\section{The Relative Zeta Factor and Its Explicit Ceiling}
\label{an:section}

We bound the relative zeta value needed in the geometric construction.
Each zeta function is normalized by the absolute degree of its field.
With this normalization, the relative degree contributes no additional
factor of two to the final bound.

Let $K$ be a field in Theorem~\ref{tw:field-family}, let $c_0$ be its
actual complex conjugation, and put
\[
 F=K^{\langle c_0\rangle},\qquad [K:\mathbb Q]=2d,\qquad
 [F:\mathbb Q]=d=b+2c,\qquad \theta=c/d.
\]
The extension $K/F$ is quadratic and unramified at all finite primes.
The field $K$ is totally imaginary and $b\geq1$. Its root discriminant, equal to
that of $F$, is
\[
 \lambda=2^{9/4}\sqrt{15015},\qquad
 \ell=\log\lambda=\frac94\log2+\frac12\log15015.
\]
Write $L_F(s)=\zeta_K(s)/\zeta_F(s)$. We will prove the following bound.

\begin{theorem}\label{an:ceiling}
For all sufficiently large degrees in the constructed field family,
\[
 \frac{\log L_F(1)}d<0.042161819.
\]
No generalized Riemann hypothesis, effective lower bound for the first
such degree, or asymptotic class-number formula is required.
\end{theorem}

\subsection{Finite-Field Comparison and Passage to One}

For a number field $A$, denote its class number, regulator, and number
of roots of unity by $h_A,R_A,w_A$. The analytic class-number formula
gives the positive value
\begin{equation}\label{an:class-formula}
 L_F(1)=2^c\pi^{d-c}\lambda^{-d/2}
       \frac{h_KR_Kw_F}{h_FR_Fw_K}>0.
\end{equation}
Indeed, the residues of $\zeta_A$ are
$2^{r_1(A)}(2\pi)^{r_2(A)}h_AR_A/(w_A\sqrt{|\Delta_A|})$,
and finite unramifiedness gives $|\Delta_K|=|\Delta_F|^2$.
We bound the factor in~\eqref{an:class-formula}. The relative unit and
ideal-class normalizations enter later in the geometric argument.

At a prime of $F$ of norm $q$, quadratic splitting gives respectively
\[
 \zeta_{K,v}(s)=(1-q^{-s})^{-2}
 \quad\hbox{or}\quad
 \zeta_{K,v}(s)=(1-q^{-2s})^{-1}.
\]
Both are at most $\zeta_{F,v}(s)^2$ for real $s>1$. Consequently
\begin{equation}\label{an:relative-euler}
 \frac{\log L_F(s)}d
 \leq\frac{\log\zeta_K(s)}{2d}.
\end{equation}
For a rational prime $p$ with absolute ramification index $e$ and
residue degree $f$ in a Galois field, its contribution to the normalized
logarithm is
\begin{equation}\label{an:local-phi}
 \Phi_{p,s}(e,f)=\frac{-\log(1-p^{-fs})}{ef}.
\end{equation}
Multiplying either $e$ or $f$ by a positive integer cannot increase
this expression. Hence, in an extension of finite Galois fields, the
normalized zeta function of the larger field is at most that of the
smaller field, prime by prime.

Set $\Gamma_{\mathbb R}(s)=\pi^{-s/2}\Gamma(s/2)$ and
$\Gamma_{\mathbb C}(s)=2(2\pi)^{-s}\Gamma(s)$. The nontrivial
quadratic Hecke character of $K/F$ is unramified at finite places and
odd at each real place of $F$. The Hecke continuation and functional
equation theorem \cite{Hecke1920,Tate1967} gives, up to a
positive constant independent of $s$, the completion
\begin{equation}\label{an:completion}
 \mathcal L_F(s)=\lambda^{ds/2}\pi^{-bs/2}(2\pi)^{-cs}
       \Gamma((s+1)/2)^b\Gamma(s)^cL_F(s).
\end{equation}
It is entire of order one. Taking the quotient of the completed
Dedekind zeta functional equations shows that its root number is $+1$.

The zero pairing in the functional equation lets us compare this
completion at one with its values to the right of one.
\begin{lemma}\label{an:monotonicity}
The function $\mathcal L_F(s)$ is nondecreasing on the real interval
$[1,\infty)$.
\end{lemma}
\begin{proof}
The Euler product and functional equation place every zero in
$0\leq\operatorname{Re}s\leq1$. Center the function at $1/2$ and
write $t=s-1/2$. Its zeros occur in pairs $z,-z$ and in conjugate
pairs. Since its order is one, $\sum_{z\ne0}|z|^{-2}<\infty$.
It follows that the paired canonical product and its logarithmic
derivative converge locally off the zeros. Evenness makes the remaining linear exponential
factor constant. A zero of multiplicity $m$ at the center contributes
$m/t$ to the logarithmic derivative. A pair with $z=a+iv$ contributes
the real part
\[
 \operatorname{Re}\frac{2t}{t^2-z^2}
 =\frac{2t(t^2-a^2+v^2)}
 {(t^2-a^2+v^2)^2+4a^2v^2}\geq0\qquad(t\geq1/2),
\]
because $|a|\leq1/2$. There are no zeros on $[1,\infty)$:
this follows from the Euler product above one and
\eqref{an:class-formula} at one. The completion is positive there.
Its logarithmic derivative is nonnegative, proving the lemma.
\end{proof}

For $\varepsilon>0$, division of the two completion factors gives
\begin{align}
 \frac{\log L_F(1)}d
 &\leq\frac{\log\zeta_K(1+\varepsilon)}{[K:\mathbb Q]}
           +G(\varepsilon,\theta),\label{an:finite-gamma}\\
 G(\varepsilon,\theta)
 &=\frac\varepsilon2(\ell-\log\pi)-\varepsilon\theta\log2
 +(1-2\theta)\log\Gamma(1+\varepsilon/2)
 +\theta\log\Gamma(1+\varepsilon).
 \label{an:gamma}
\end{align}
In particular $G(\varepsilon,\theta)\to0$ uniformly for
$0\leq\theta\leq1/2$ as $\varepsilon\downarrow0$.

\subsection{An Unconditional Asymptotic Prime Budget}

The basic inequality for limiting prime counts improves a finite Euler
estimate at a fixed abscissa to the following bound at one. Only this
inequality is needed, without a Brauer--Siegel equality.

\begin{proposition}\label{an:prime-budget}
Suppose that, at $\sigma=1+\varepsilon>1$, every field in the family
satisfies
\[
 \frac{\log\zeta_K(\sigma)}{[K:\mathbb Q]}\leq Y_K(\sigma)
 \leq Y_*(\sigma).
\]
Let $e_p^0,f_p^0$ be uniform lower bounds for its absolute local indices,
and suppose
\[
 R\geq\sum_p\frac{\log p}{e_p^0(p^{2f_p^0}-1)}.
\]
Then
\begin{equation}\label{an:budget-conclusion}
 \limsup_{d\to\infty}\frac{\log L_F(1)}d
 \leq Y_*(1+\varepsilon)
 +\varepsilon\left(\frac{\ell-\gamma_E-\log(4\pi)}4+R\right).
\end{equation}
Every strictly larger constant is an eventual uniform ceiling.
\end{proposition}
\begin{proof}
Take any sequence with $n=[K:\mathbb Q]=2d\to\infty$. A diagonal
subsequence makes each normalized prime count converge:
$\beta_q=\lim N_q(K)/n$, where $q$ ranges over prime powers.
The fixed root discriminant makes this an asymptotically exact family.
The unconditional Basic Inequality$'$ of
\cite[Proposition 3.1]{TsfasmanVladut2002} is
\begin{equation}\label{an:tv-budget}
 \sum_q\beta_q w(q)\leq B,\qquad
 w(q)=2\log q\sum_{m\geq1}\frac1{q^m+1},\qquad
 B=\frac{\ell-\gamma_E-\log(4\pi)}2.
\end{equation}
To check the normalization, their genus is $g=n\ell/2$, their finite
invariants are $\phi_q=2\beta_q/\ell$, and total imaginary signature
gives $\phi_{\mathbb R}=0$, $\phi_{\mathbb C}=1/\ell$.
The proposition is unconditional, with archimedean constant
$\gamma_E+\log(4\pi)$.

Put $a_s(q)=-\log(1-q^{-s})$ and
$Z_\infty(s)=\sum_q\beta_q a_s(q)$. The series converges at one,
since $a_1(q)/w(q)\leq(q+1)/(2(q-1)\log q)$ is bounded for $q\geq2$.
For every fixed $\tau>1$, dominated convergence in the normalized
Euler products gives
\[
 \lim\frac{\log\zeta_K(\tau)}n=Z_\infty(\tau).
\]
One can dominate prime by prime by $-\log(1-p^{-\tau})$.
Equation~\eqref{an:finite-gamma}, followed by $\tau\downarrow1$,
then gives $\limsup d^{-1}\log L_F(1)\leq Z_\infty(1)$.
The uniform vanishing of $G$ and monotone convergence justify this
order of limits. The finite-field Euler products are used only to the
right of one.

For every prime power $q$,
\begin{align*}
 a_1(q)-a_{1+\varepsilon}(q)&\leq
                \frac{\varepsilon\log q}{q-1},\\
 \frac{\log q}{q-1}-\frac{w(q)}2
 &=\log q\sum_{m\geq1}\frac1{q^m(q^m+1)}
 \leq\frac{\log q}{q^2-1}.
\end{align*}
For a Galois field, the normalized contribution above $p$ to the last
weighted sum is $\log p/[e_p(p^{2f_p}-1)]$. The local floors and
Fatou's lemma bound its limiting sum by $R$. Combining this with
\eqref{an:tv-budget} proves
$Z_\infty(1)\leq Z_\infty(1+\varepsilon)+\varepsilon(B/2+R)$.
This proves the assertion on every such subsequence. If a strictly
larger constant were not an eventual uniform bound, a sequence of
violating fields would have a diagonal subsequence of the kind just
considered, giving a contradiction.
\end{proof}

\subsection{The Complete Finite-Field Euler Factorization}
\label{an:factorization}

Let $N$ be the degree-$16384$ retained field of
Proposition~\ref{ar:central-fields}, and set
$S'=\{2,3,5,7,11,13,17\}$. A superscript $S'$ removes these Euler
factors. Proposition~\ref{ar:hecke-data} gives all irreducible
representations of $\operatorname{Gal}(N/\mathbb Q)$ and realizes
each nonlinear factor as an induced nontrivial finite-order Hecke
character. There are $128$ linear, $256$ quadratic, $456$ quartic,
and $124$ octic factors. Among the quartics, eight have pure gamma
parity, and every other nonlinear factor is mixed. The regular
representation identity is
\begin{equation}\label{an:regular-factorization}
 \zeta_N(s)=\prod_\rho L(s,\rho)^{\dim\rho},\qquad
 128+256\cdot4+456\cdot16+124\cdot64=16384.
\end{equation}
For real $s>1$ the factors are nonzero. Real factors are positive:
their Euler products tend to one as $s\to\infty$ and cannot cross
zero. Conjugate factors have product $|L(s,\rho)|^2$. Thus
\begin{equation}\label{an:Y}
 Y(s):=\frac{\log\zeta_N^{S'}(s)}{16384}
 =\frac1{16384}\sum_\rho\dim\rho\,
          \log|L^{S'}(s,\rho)|.
\end{equation}
Each conjugate pair is counted twice in this sum, with the same
representation degree. An upper bound for each modulus may be
substituted after this identity, even when some logarithms are negative.

The exact arithmetic input for one factor consists of its conductor
$Q$, gamma parities, root number when proved to be one, an inducing
radicand, genus twist, and the exceptional denominators
\[
 D_{\rho,p}(T)=\det(1-T\rho(\operatorname{Frob}_p)
                         \mid V_\rho^{I_p}).
\]
Their construction, including the cancellation of quartic twists at
5 and 13, is given in Proposition~\ref{ar:hecke-data}. The
analytic evaluations below use the representations constructed there,
including their conductors and local characters. Formal compatibility
among a set of Euler polynomials would not identify these functions.

\subsection{Exact Coefficients and Their Enumeration}
\label{an:coefficients}

Write $L(s,\rho)=\sum_{n\geq1}a_n n^{-s}$ and $r=\dim\rho$.
At every prime, the reciprocal roots of $D_{\rho,p}$ have modulus
one, and there are at most $r$ of them. Hence
\begin{equation}\label{an:divisor-bound}
 |a_n|\leq d_r(n),\qquad
 \sum_{n\geq1}d_r(n)n^{-\beta}=\zeta(\beta)^r
 \quad(\beta>1).
\end{equation}
The finite-image representation also makes every coefficient used here
an exact Gaussian integer.

To construct the coefficients outside $S'$, let
$v_p\in\mathbb F_2^7$ be genus Frobenius, and let $B_q$ be the polar
form of the corresponding central character. Outside its radical,
Frobenius has eigenvalues in opposite pairs and square $\kappa I$,
where $\kappa\in\{1,-1\}$. Its denominator is
$(1-\kappa T^2)^{r/2}$, before the explicitly supplied quartic twist.
The coefficients at exponents $2j$ are
$\binom{j+r/2-1}{r/2-1}\kappa^j$, and odd exponents vanish.
Inside the radical, Frobenius is scalar, and its denominator is
$(1-\omega T)^r$, where $\omega\in\{1,-1,i,-i\}$ is determined
by the inducing field and twist. Its coefficients are
$\binom{j+r-1}{r-1}\omega^j$.

At a scalar prime the inducing multiquadratic base splits completely.
For the positive quadratic induction, the scalar sign is obtained by
evaluating the radicand at these embeddings modulo $p$ and taking
quadratic residue symbols. These symbols agree
because the represented Frobenius is scalar. There is an exact
procedure even at a prime dividing the norm of the displayed radicand.
For an integral radicand $\eta$, let $m=v_p(N\eta)$. Each embedding
valuation is between zero and $m$. Lift the base square roots to
modulus $p^{m+1}$ by Hensel's lemma. This determines every embedding
valuation and normalized nonzero unit residue. The valuations are
even, by unramifiedness outside $S'$. The normalized quadratic symbols
give the required signs. Rational square normalization does not change
the extension. This procedure treats every norm prime, rather than
assuming that a finite list of ordinary exceptions is exhaustive.

The genus and twist characters are evaluated exactly by Kronecker
symbols and by discrete logarithms in $(\mathbb Z/5\mathbb Z)^\times$
or $(\mathbb Z/13\mathbb Z)^\times$, with $\xi_p(2)=i$.
The genus characters have common period $120120$, so lookup tables may
use this period or any multiple of it. At $p\in S'$, the
verified polynomial $D_{\rho,p}$ supplies the recurrence
\[
 a_{p^j}=-\sum_{k=1}^{\min(j,r)}[T^k]D_{\rho,p}(T)\,a_{p^{j-k}},
 \qquad a_1=1.
\]
Multiplicativity then determines all coefficients.

To enumerate the nonzero coefficients through the cutoff $N_0$, first
list all primes up to $N_0$ by a segmented sieve using every sieving
prime through $\lfloor\sqrt{N_0}\rfloor$. Assign an ordinary nonscalar prime the
first possible support location $p^2$, a scalar prime the location $p$,
and an exceptional prime the safe lower bound $p$. Sort primes by these
locations. Recursion appends only later primes in that order, considering
every admissible prime power and continuing across intermediate zero
coefficients. Unique factorization then visits each supported integer
exactly once. If a first support location exceeds the remaining quotient,
that prime and every later prime can be skipped. These observations prove that
the enumeration is complete.

Signed and absolute moments use arbitrary-precision integers. Fixed-width
coefficient arithmetic is checked before conversion, and intermediate
products are formed in a wider type. An overflow aborts the calculation.
Modular multiplication uses a wide product or a checked Montgomery
reduction as appropriate. The coefficients, support counts, and moments are computed exactly.

\subsection{The Linear Factors}

The $128$ linear factors are the primitive quadratic Dirichlet functions
of the genus field
$E=\mathbb Q(i,\sqrt2,\sqrt3,\sqrt5,\sqrt7,\sqrt{11},\sqrt{13})$,
including the principal factor $\zeta(s)$. For a signed squarefree
product $D_0$ of these radicands, its fundamental discriminant is
$D=D_0$ if $D_0\equiv1\pmod4$, and $D=4D_0$ otherwise.
The character is the Kronecker symbol $(D/\cdot)$.
The implementation checks its primitive conductor and its values on
explicit generators of every prime-power component of the unit group.
These values identify the character independently of its library index.

The finite identity
\begin{equation}\label{an:dirichlet-hurwitz}
 L(s,\chi)=q^{-s}\sum_{a=1}^{q}\chi(a)\zeta(s,a/q)
\end{equation}
reduces evaluation to Hurwitz zeta functions. Euler--Maclaurin summation
gives the explicit enclosure
\begin{align*}
 \zeta(s,x)={}&\sum_{k=0}^{M-1}(k+x)^{-s}
 +\frac{(M+x)^{1-s}}{s-1}+\frac12(M+x)^{-s}\\
 &+\sum_{j=1}^{J}\frac{B_{2j}}{(2j)!}(s)_{2j-1}
                     (M+x)^{1-s-2j}+\mathcal R,\\
 |\mathcal R|\leq{}&
 \frac{2\zeta(2J)(s)_{2J}}{(2\pi)^{2J}(s+2J-1)}
                       (M+x)^{1-s-2J},
 \qquad s>1, x>0.
\end{align*}
The remainder follows from the periodic Bernoulli bound
$|\widetilde B_{2J}(t)|\leq2(2J)!\zeta(2J)/(2\pi)^{2J}$.
For rigorous high-precision evaluation by Euler--Maclaurin summation,
see \cite{Johansson2015Hurwitz}. The supplied evaluation uses Arb/FLINT
ball arithmetic \cite{Johansson2017}, retaining the endpoints of each
real value.
Removing $S'$ multiplies a factor by
$\prod_{p\in S'}(1-\chi(p)p^{-s})$.

\subsection{An Approximate Functional Equation for Mixed Factors}
\label{an:mixed-afe}

Numerical evaluation from gamma factors and a functional equation is
developed in \cite{Dokchitser2004}. We use this approach with explicit
kernel and tail bounds for the factors constructed above.
For a mixed factor of degree $2m$, where $m=1,2,4$, its completion is
\[
 \Lambda(s)=Q^{s/2}\Gamma_{\mathbb R}(s)^m
                    \Gamma_{\mathbb R}(s+1)^m L(s)
 =2^m Q^{s/2}(2\pi)^{-ms}\Gamma(s)^mL(s).
\]
It is entire and satisfies
$\Lambda(s)=\varepsilon_\rho\overline{\Lambda(1-\overline s)}$,
with $|\varepsilon_\rho|=1$. For a real quadratic induction the
Dedekind zeta quotient gives $\varepsilon_\rho=1$.

Let $k_m(x)$ have Mellin transform $\Gamma(z)^m$, and put
\[
 H_{m,a}(x)=\int_1^\infty u^{a-1}k_m(xu)\,du,\qquad
 t=\frac{(2\pi)^m}{\sqrt Q},\qquad
 F_m=\frac{t^s}{\Gamma(s)^m}.
\]
These kernels are positive. More explicitly,
\begin{equation}\label{an:positive-kernel}
 k_m(x)=\int_{(0,\infty)^{m-1}}
 \exp\left(-u_1-\cdots-u_{m-1}-\frac{x}{u_1\cdots u_{m-1}}\right)
                  \frac{du_1\cdots du_{m-1}}{u_1\cdots u_{m-1}},
\end{equation}
with $k_1(x)=e^{-x}$. Taking the Mellin transform under this positive
integral proves the asserted transform and also complete monotonicity
in $x$. In particular $k_2(x)=2K_0(2\sqrt x)$
\cite[10.32.10, 10.43.19]{DLMF}.

Splitting the Mellin integral at one gives a common approximate
functional equation for the mixed factors.
\begin{lemma}\label{an:balanced-afe}
For real $s>1$ and every nonlinear mixed factor,
\begin{equation}\label{an:mixed-formula}
 L(s)=F_m\sum_{n\geq1}a_nH_{m,s}(tn)
 +\varepsilon_\rho F_m\sum_{n\geq1}\overline{a_n}
                                      H_{m,1-s}(tn).
\end{equation}
\end{lemma}
\begin{proof}
The function $\Theta(u)=\sum a_n k_m(tnu)$ has Mellin transform
$\Gamma(z)^m t^{-z}L(z)=2^{-m}\Lambda(z)$. Mellin inversion and
the completed functional equation give
$\Theta(u)=\varepsilon_\rho u^{-1}\overline{\Theta(1/u)}$.
The contour shift has no residues because the completion is entire,
and the gamma factors give the needed vertical decay. Splitting the Mellin
integral at one and substituting $u\mapsto1/u$ below one yields
\eqref{an:mixed-formula}. The positive integral representation and
\eqref{an:divisor-bound} justify the sums in either exponentially
decaying half.
\end{proof}

For degree two this uses
$H_{1,a}(x)=\int_1^\infty u^{a-1}e^{-xu}\,du$.
For the mixed quartic implementation it is convenient to write
\[
 I_b(y)=\int_1^\infty u^bK_0(yu)\,du,\qquad
 H_{2,a}(x)=4I_{2a-1}(2\sqrt x).
\]
Its prefactor is $4t^s/\Gamma(s)^2$ when the stored kernel
is $I_{2s-1}+I_{1-2s}$. This identity accounts for the factor four
in the quartic evaluation. For octics the stored kernel is
$H_{4,a}$ itself, with prefactor $t^s/\Gamma(s)^4$.

\subsection{Tail Bounds and Absolute Coefficient Prefixes}
\label{an:tails}

For $\beta\geq\max(s,1-s)$,
$H_{m,a}(x)\leq H_{m,\beta}(x)$ when $a\leq\beta$, and
\[
 x^\beta H_{m,\beta}(x)=\int_x^\infty v^{\beta-1}k_m(v)\,dv
\]
is decreasing. If $M_{N_0}(\beta)$ bounds
$\sum_{n>N_0}|a_n|n^{-\beta}$, either omitted half of
\eqref{an:mixed-formula} is at most
\begin{equation}\label{an:general-tail}
 F_m (N_0+1)^\beta H_{m,\beta}(t(N_0+1))M_{N_0}(\beta).
\end{equation}
One may always use $M_{N_0}(\beta)=\zeta(\beta)^{2m}$.

The quadratic and quartic evaluations also use elementary envelopes.
For the exponents in question, with $1<s<1.01$,
\begin{align}
 H_{1,s}(x),H_{1,1-s}(x)&\leq
 V(x):=e^{-x}(1+1/x)/x,\label{an:exp-envelope}\\
 I_{2s-1}(y),I_{1-2s}(y)&\leq
 U(y):=\sqrt{\pi/2}\,e^{-y}(1+1/y)y^{-3/2}.
 \label{an:bessel-envelope}
\end{align}
The first follows by replacing $u^{s-1}$ and $u^{-s}$ by $u$.
For the second use $K_0(yu)\leq\sqrt{\pi/(2yu)}e^{-yu}$ and
$u^{b-1/2}\leq u$ for $b\leq3/2$.
The functions $n^\beta V(tn)$ and
$n^\beta U(2\sqrt{tn})$ decrease respectively when
$tn>\beta-1$ and $2\sqrt{tn}>2\beta-3/2$.
Thus substituting these envelopes and the divisor majorant gives
valid tails even when the chosen $\beta$ exceeds an exponent covered
by the envelope. The quartic prefactor includes the factor four from the conversion
between the kernels. When we sum both halves, we include both tail
allowances.

For the octics a sharper absolute series is available. Let
$W=\operatorname{rad}B_q$. Then $[E^W:\mathbb Q]=64$. At an ordinary
prime, absolute coefficients have local series $(1-T)^{-8}$ if
$v_p\in W$, and $(1-T^2)^{-4}$ otherwise. These are precisely the
local factors of $\zeta_{E^W}^{S'}(\beta)^{1/8}$.
For every exceptional denominator in the supplied arithmetic table,
an exact Gaussian-polynomial comparison supplies nonnegative integers
$u_p,v_p$ and a fourth root of unity $\omega_p$ such that
\begin{equation}\label{an:exceptional-phase}
 D_{\rho,p}(T)=(1-\omega_pT)^{u_p}
                          (1-\omega_p^2T^2)^{v_p}.
\end{equation}
The check enumerates $u_p,v_p\geq0$ with $u_p+2v_p\leq8$ and
the four possible $\omega_p$, comparing all polynomial coefficients
exactly. After the rotation $T\mapsto\omega_p^{-1}T$, the reciprocal
has nonnegative coefficients. Consequently the full absolute series is
exactly
\begin{equation}\label{an:absolute-series}
 M_\rho(\beta):=\sum_{n\geq1}|a_n|n^{-\beta}
 =\zeta_{E^W}^{S'}(\beta)^{1/8}
 \prod_{p\in S'}(1-p^{-\beta})^{-u_p}
                       (1-p^{-2\beta})^{-v_p}.
\end{equation}
The same rotations show that each nonzero coefficient lies on a
Gaussian coordinate axis, so $|a_n|=|\operatorname{Re}a_n|+
|\operatorname{Im}a_n|$ exactly. The genus subfield zeta value is
evaluated using its $64$ quadratic characters by
\eqref{an:dirichlet-hurwitz}.

We subtract a rigorous lower bound for the finite absolute prefix.
On a bin with integer midpoint $m_0$ and radius $r_0$, record exact
moments $A_j=\sum|a_n|(n-m_0)^j$, $0\leq j\leq J=20$.
Let $h=\lceil\beta\rceil$ and $\rho_0=r_0/m_0<1/16$.
The degree-$J$ binomial expansion of $n^{-\beta}$ has error at most
\begin{equation}\label{an:absolute-prefix-error}
 A_0m_0^{-\beta}\binom{J+h}{h-1}
                    \frac{\rho_0^{J+1}}{(1-\rho_0)^h}.
\end{equation}
Indeed $(\beta)_j/j!\leq\binom{j+h-1}{h-1}$, and after the first
omitted coefficient their ratios are bounded by the coefficients of
$(1-\rho_0)^{-h}$. Sum the interval polynomials and subtract these
errors to obtain $P_{N_0}(\beta)\leq\sum_{n\leq N_0}|a_n|n^{-\beta}$.
Then $M_\rho(\beta)-P_{N_0}(\beta)$ is an admissible
$M_{N_0}(\beta)$ in~\eqref{an:general-tail}.
The octic evaluation selects the smallest valid enclosure among
$\beta=5/4,3/2,2,5/2,3$. Each candidate is an upper bound, so this selection preserves the
inequality.

\subsection{The Eight Pure Quartic Factors}
\label{an:pure}

The pure factors in Proposition~\ref{ar:hecke-data} are real quadratic
Hecke inductions with root number one, conductor $Q=240240^2$, and
gamma vector $(a,a,a,a)$, $a\in\{0,1\}$, four of each parity.
Their completion is $Q^{s/2}\Gamma_{\mathbb R}(s+a)^4L(s)$.
Put
\[
 t=\pi^2/\sqrt Q,\quad b_0=(s+a)/2,\quad b_1=(1-s+a)/2,
 \quad F=\frac{t^{s+a}}{\Gamma(b_0)^4},\quad K_b=H_{4,b}.
\]
Then the exact functional equation is
\begin{equation}\label{an:pure-afe}
 L(s)=F\sum_{n\geq1}a_n n^a
       \bigl(K_{b_0}(t^2n^2)+K_{b_1}(t^2n^2)\bigr).
\end{equation}
To verify the normalization, the Mellin transform of
$\sum a_n n^a k_4(t^2n^2u)$ is
\[
 \Gamma(z)^4t^{-2z}L(2z-a)=Q^{a/2}\Lambda(2z-a).
\]
Reflection is about $z=(a+1/2)/2$. Splitting at one gives
\eqref{an:pure-afe}, with no extra factor two.
For $\beta\geq\max(s,1-s)$, put $B=(\beta+a)/2$. Either omitted
side is bounded by
\begin{equation}\label{an:pure-tail}
 F(N_0+1)^{\beta+a}K_B(t^2(N_0+1)^2)\zeta(\beta)^4.
\end{equation}
This follows from the same decreasing-integral argument as
\eqref{an:general-tail}. The exact cutoff here is
$N_0=8\sqrt Q=1921920$.

\subsection{Kernel and Moment Evaluation}
\label{an:kernel-enclosures}

We evaluate the kernels on finite grids and bound the interpolation
errors. For the four-gamma kernel, let
$H_n^{(j)}=\sum_{k=1}^n k^{-j}$ and
\[
 L_n=4(H_n^{(1)}-\gamma_E)-\log x,\quad
 P_{3,n}=L_n^3/6+2L_n(\zeta(2)+H_n^{(2)})
                     +4(H_n^{(3)}-\zeta(3))/3,
\]
\[
 P_{2,n}=L_n^2/2+2(\zeta(2)+H_n^{(2)}).
\]
Residues at the order-four poles of $\Gamma(z)^4$ give the convergent
series
\begin{align}
 k_4(x)&=\sum_{n\geq0}\frac{x^n}{(n!)^4}P_{3,n},
 \label{an:k4-series}\\
 K_a(x)&=\Gamma(a)^4x^{-a}-
 \sum_{n\geq0}\frac{x^n}{(n!)^4}
 \left(\frac{P_{3,n}}{n+a}+\frac{P_{2,n}}{(n+a)^2}
       +\frac{L_n}{(n+a)^3}+\frac1{(n+a)^4}\right).
 \label{an:K4-series}
\end{align}
The second identity initially follows for $a>0$ and extends to
$-1/100<a<0$ by analytic continuation. None of the exponents evaluated
here is zero or a negative integer, so no colliding-pole formula is
needed.

To bound the tails of these series, for $n\geq2$ set
\[
 B_n=|\log x|+4(1+\log n+\gamma_E),\quad
 U_{3,n}=B_n^3/6+4\zeta(2)B_n+4\zeta(3)/3,
 \quad U_{2,n}=B_n^2/2+4\zeta(2).
\]
They bound $|L_n|,|P_{3,n}|,|P_{2,n}|$ in the corresponding order.
Also $B_{n+1}\leq(1+1/n)B_n$, so
\[
 r_n=\frac{x}{(n+1)^4}(1+1/n)^3
\]
is an upper bound for the subsequent term ratios of the majorants.
Once $r_n<1$, the entire omitted tail of~\eqref{an:k4-series}
is at most
\[
 \frac{x^n}{(n!)^4}\frac{U_{3,n}}{1-r_n}.
\]
For~\eqref{an:K4-series}, provided $a>-1/100$, it is at most
\[
 \frac{x^n}{(n!)^4}\frac1{1-r_n}
 \left(\frac{U_{3,n}}{n-1}+\frac{U_{2,n}}{(n-1)^2}
       +\frac{B_n}{(n-1)^3}+\frac1{(n-1)^4}\right).
\]
All these expressions are evaluated outward. The first four Taylor
coefficients of $k_4$ are obtained by the same argument on
$|z-x|\leq x/2$, replacing $x$ by $3x/2$ in the numerator bound
and $|\log x|$ by $|\log x|+\log2+\pi/2$.
Cauchy's estimate bounds the error in the $j^{th}$ coefficient by
the uniform tail times $(2/x)^j$.

Higher Taylor coefficients follow from
\begin{equation}\label{an:kernel-odes}
 (x\partial_x)^4k_4=xk_4,\qquad xK_a'+aK_a=-k_4.
\end{equation}
These identities follow directly from the Mellin transform and the
defining integral. They give finite recurrences using rational
operations and the enclosed initial values. The two-gamma kernel
uses $K(x)=K_0(2\sqrt x)$, $J_b(x)=I_b(2\sqrt x)$, with
\[
 xK''+K'-K=0,\qquad 2xJ_b'+(b+1)J_b=-K.
\]
If $k_j,j_j$ denote Taylor coefficients at $x_0$, the recurrences are
\[
 k_{j+2}=\frac{k_j-(j+1)^2k_{j+1}}{x_0(j+1)(j+2)},\qquad
 j_{j+1}=\frac{-k_j-(2j+b+1)j_j}{2x_0(j+1)}.
\]
The initial values are $K_0(2\sqrt{x_0})$ and
$-K_1(2\sqrt{x_0})/\sqrt{x_0}$. The one-gamma case uses the elementary
exponential and $xH_{1,a}'+aH_{1,a}=-e^{-x}$.

The grids have consecutive endpoints $x_0,31x_0/32$.
For any positive Laplace kernel $H$ in use, its holomorphic extension
obeys $|H(z)|\leq H(x_0/2)$ on $|z-x_0|\leq x_0/2$.
Hence a degree-$J$ Taylor polynomial has error at most
\begin{equation}\label{an:taylor-error}
 H(x_0/2)\frac{16^{-J-1}}{1-1/16}
\end{equation}
throughout the bin. The backward Taylor remainder is nonnegative by
complete monotonicity, allowing propagation of the exponential and
Bessel kernels from a certified large-argument interval.
Their starting values are bounded by~\eqref{an:exp-envelope} and
\eqref{an:bessel-envelope}. The four-gamma kernel instead uses the
residue series at each center. Grid endpoints, their integer floors,
and complete coverage from $1$ to the exact coefficient cutoff are
checked with outward intervals. The calculation is rejected if an
integer floor is unresolved.

For a linear kernel argument $x=tn$, translating its degree-$20$
Taylor polynomial at $x_0$ to the integer midpoint $m_0$ gives
\[
 [ (n-m_0)^j ]\,P(tn)
 =t^j\sum_{h=j}^{20}[ (x-x_0)^h ]P(x)
                 \binom hj(tm_0-x_0)^{h-j}.
\]
Contract this polynomial against the exact real and imaginary moments.
Multiplying~\eqref{an:taylor-error} by the sum of absolute coefficients
in the bin bounds its contribution to the error.
For the pure quartics the substitution is quadratic:
\[
 t^2n^2-x_0=(t^2m_0^2-x_0)+2t^2m_0(n-m_0)+t^2(n-m_0)^2.
\]
A degree-$20$ kernel polynomial thus needs moments through degree
$40$. These moments and their absolute masses include $n^a$ exactly
once. The uniform remainder still has degree $20$ in the kernel
variable, so~\eqref{an:taylor-error} applies before composition.

Suppose $A,B$ are the finite primary and conjugated dual sums, and
$E_A,E_B$ bound their interpolation and infinite-tail errors. A known
root number one gives $|L(s)|\leq|A+B|+E_A+E_B$. With an unknown
phase one always has $|L(s)|\leq|A|+|B|+E_A+E_B$.
Every modulus, gamma value, logarithm, and conversion to an endpoint
retains its rounding enclosure.

\subsection{Recovering the Mixed Quadratic and Quartic Phase}
\label{an:two-cutoffs}

For the complex quadratic and mixed quartic factors, a second Mellin
cutoff can sharpen this bound by enclosing the root number. Split the
Mellin integral at $y>0$ instead of one. The two exact terms are
\[
 A_y=F_m y^s\sum a_nH_{m,s}(tyn),\qquad
 B_y=F_m y^{s-1}\sum\overline{a_n}H_{m,1-s}(tn/y),
 \qquad L(s)=A_y+\varepsilon_\rho B_y.
\]
Take $y=(31/32)^{22}$, so shifted kernels reuse the same exact moment
bins. Any terminal part outside the stored kernel range is put into
the tail with its actual first omitted integer.

For finite approximations, write $U=A_y-A_1$, $V=B_1-B_y$ and let
$\eta,\Delta$ bound their respective errors. If $|V|>\Delta$, then
\begin{equation}\label{an:phase-disk}
 \varphi=U/V,\qquad
 |\varepsilon_\rho-\varphi|
 \leq r_\varphi:=\frac{\eta+|\varphi|\Delta}{|V|-\Delta}.
\end{equation}
This follows by writing the exact quotient as
$(U+u)/(V+v)$, with $|u|\leq\eta$, $|v|\leq\Delta$.
If each balanced side has error at most $E_1$, the resulting bound is
\begin{equation}\label{an:phase-upper}
 |L(s)|\leq|A_1+\varphi B_1|+2E_1+|B_1|r_\varphi.
\end{equation}
The coefficient $2E_1$ uses the known identity
$|\varepsilon_\rho|=1$, not an assumption that $|\varphi|=1$.
All distances from zero and radii are enclosed outward. If the
denominator cannot be separated from zero, the phase-free bound is
retained. Complex octic factors use the phase-free bound directly.

\subsection{The Certified Factor Sums}
\label{an:factor-sums}

The arithmetic table and exact moment files in the supplementary
certificate determine every factor just described. The abscissa is
the exact rational
\[
 \sigma=\frac{12001}{12000}.
\]
For each row, the evaluator applies the appropriate AFE, includes
interpolation and both infinite tails, and removes the exceptional
factors by
\[
 \log|L^{S'}(\sigma,\rho)|
 \leq\log U_\rho+
       \sum_{p\in S'}\log|D_{\rho,p}(p^{-\sigma})|,
\]
where $U_\rho$ is its proved modulus upper bound.
The octic cutoffs are
$N_0=\min(3\cdot10^9,\lfloor\sqrt Q/40\rfloor)$.
The exact integers for the quadratic and mixed quartic rows are
specified in their moment records and individually checked when
constructing the bins and tails. The pure quartic cutoff was given
in Section~\ref{an:pure}.

The following table gives rational upper bounds for the sums of
logarithms. The number of factors includes both members of each
conjugate pair. An interval enclosing an AFE upper bound encloses that
bound, and need not give a two-sided enclosure of comparable width for
the $L$-value itself.
\begin{center}
\begin{tabular}{lrr}
\toprule
Type & Number & Upper bound for $\sum\log|L^{S'}(\sigma,\rho)|$\\
\midrule
Linear &128& $7.337578130846530221$\\
Mixed quadratic &256& $0.400886873183882259$\\
Mixed quartic &448& $-0.046453348154273950$\\
Pure quartic &8& $0.058145052434373409$\\
Mixed octic &124& $-0.111176963043772737$\\
\bottomrule
\end{tabular}
\end{center}
Exact addition of their dimension-weighted upper endpoints in
\eqref{an:Y} gives
\begin{equation}\label{an:Y-numerical}
 Y(\sigma)\leq0.00044535540710354679437255859375
             <0.000445355407103547.
\end{equation}
The coefficient algorithms, kernel bounds, and tail estimates above
specify how to verify these inequalities. The supplementary program
checks the exact row identities and partitions before performing the
interval evaluations. Its moment mode regenerates the newest 64 sector
moment files from the coefficient algorithms. Earlier moment records
remain finite inputs, and the data manifest checks their identities.
% CHECK: The current replay does not regenerate all earlier moments or actual-field data.
% The required independent verification scope is recorded in evidence/body-review-arithmetic-20260927.md.

\subsection{Finite Corrections beyond the Retained Field}
\label{an:census}

To compare $N$ with the whole tower, restore $S'$ using the actual
local indices of $K$. Its exact contribution is
\begin{align}
 B_D(\sigma)={}&-\frac{\log(1-2^{-4\sigma})}{32}
 -\sum_{p=3,5}\frac{\log(1-p^{-2\sigma})}{4}\\
 &-\sum_{p=7,11,13}\frac{\log(1-p^{-4\sigma})}{8}
 -\frac{\log(1-17^{-4\sigma})}{4}.
 \label{an:bad-restored}
\end{align}
This restores the deleted factors using the local indices of $K$.
The indices of $N$ at $S'$ may be different.

First consider $p\leq10^4$, $p\notin S'$, with nonzero genus vector
$v_p$. Its square is detected in the full retained degree-two quotient
if some form in the twelve-dimensional central space has $q(v_p)=1$.
If all seven completed forms instead have $q(v_p)=0$, then $f_N(p)=2$
whereas $4\mid f_K(p)$. Define $\mathcal P_4$ by precisely these
conditions. It contains $45$ primes. Equation~\eqref{an:local-phi}
gives the saving
\begin{equation}\label{an:order-four-saving}
 S_4=\frac14\sum_{p\in\mathcal P_4}
                     \log\frac{1+p^{-2\sigma}}{1-p^{-2\sigma}}
 >0.00003724741189544274807446.
\end{equation}
Primes for which the completed field already has order four are
excluded from this sum.

The second correction uses primes that split in $E$ and in $N$,
but not in the full twelve-dimensional retained central field.
Here $f_N(p)=1$ and $f_K(p)$ is even, so the saving is
\begin{equation}\label{an:central-saving}
 S_{\rm cen}=\sum_{\substack{p\leq10^{12}\\
                  p\text{ split in }N,\ p\text{ nonsplit in the full field}}}
                     \operatorname{atanh}(p^{-\sigma}).
\end{equation}
The support has zero genus vector, so it is disjoint from
$\mathcal P_4$. Both supports exclude $S'$. No saving from an
uncomputed tail is claimed.

To compute~\eqref{an:central-saving}, enumerate the genus-split primes.
Their residues are the $180$ classes modulo $120120$ satisfying
$p\equiv1\pmod8$ and being nonzero squares modulo each of
$3,5,7,11,13$. Construct them by the Chinese remainder theorem.
For each class, sieve the numbers $120120k+r$ through $X=10^{12}$
by all primes up to $\sqrt X$, beginning at the square of each
sieving prime. Every surviving number above $13$ is prime, since
a composite would have a prime divisor at most its square root.
The classes are disjoint and none contains a prime dividing the modulus.

For a surviving prime, evaluate the $17$ norm radicands of
Section~\ref{ar:catalog} using exact modular square roots. The symbol of
$x+y\sqrt A$ is independent of the choice of root: its product with
the conjugate is $Bz^2$, a nonzero square at a genus-split prime.
All catalog $B$ are supported on $2,3,5,7,11,13$ and all catalog
$z$ are $1$ or $2$. Thus no prime in this census divides $Bz$,
and no zero residue or unhandled norm prime occurs in these raw
predicates.
The seven completed and twelve full-field product predicates in
Proposition~\ref{ar:central-fields} are then exact splitting tests.
Order the completed tests first and count only the first nonzero
predicate. A hit in the first seven tests is already accounted for by
$N$. A later hit belongs to~\eqref{an:central-saving}. A prime on
which every predicate vanishes is not counted. This proves that the
computed complement contains each eligible prime exactly once.

For outward summation, use disjoint bins $l<p\leq h$, beginning at
$l=13$ and with next endpoint
$h=\min(X,\max(l+1,\lceil1001l/1000\rceil))$.
For complement primes in a bin record their exact number $n$ and
$u=\sum\lfloor10^{30}/p\rfloor$. With $\varepsilon=\sigma-1$,
their contribution is enclosed between
\begin{equation}\label{an:census-bin}
 h^{-\varepsilon}\frac{u}{10^{30}}
 \quad\hbox{and}\quad
 l^{-\varepsilon}\frac{u+n}{10^{30}}
                         \frac1{1-l^{-2}}.
\end{equation}
Indeed $x\leq\operatorname{atanh}x\leq x/(1-x^2)$ for $0<x<1$,
and $\sigma>1$. The only rounded quantities in the producer are the
explicit floors, whose full error is included in $u+n$.
All counts and reciprocal numerators use exact integers.
For the fixed-width census implementation, counts are at most $X$, and
the total reciprocal numerator is less than
$10^{30}(1+\log X)<2^{105}$. Thus its $64$-bit counts and $128$-bit
accumulators cannot overflow at this cutoff. Modular products use the
wider reduction described above.

The full sieve finds $293772662$ genus-split primes. Of these,
$293705608$ are detected by the full central field, and $291485548$
are already detected by $N$. The complement contains
$2220060$ primes. Summing~\eqref{an:census-bin} outward gives
\begin{equation}\label{an:census-value}
 0.00004201663320373563276702<S_{\rm cen}
          <0.00004201663670363397482805.
\end{equation}
An independent prime-list and modular-arithmetic calculation through
$10^7$ checks every bin, first-hit count, and reciprocal numerator.
Its complement has $31$ primes. This finite comparison checks the
implementation but does not replace the full sieve or its completeness
argument.

Primewise comparison now proves, for every retained tower field,
\begin{equation}\label{an:full-euler}
 \frac{\log\zeta_K(\sigma)}{[K:\mathbb Q]}
 \leq Y(\sigma)+B_D(\sigma)-S_4-S_{\rm cen}.
\end{equation}

\begin{remark}[Fixed-field form]\label{an:fixed-field-form}
The same primewise comparison applies to the retained field $M$ of
Proposition~\ref{ar:central-fields} itself, of degree $2^{19}$, and
gives
\[
 \frac{\log\zeta_M(\sigma)}{2^{19}}
 \leq Y(\sigma)+B_D(\sigma)-S_4-S_{\rm cen}.
\]
Since $M$ contains $N$, equation~\eqref{an:local-phi} bounds its
normalized zeta logarithm prime by prime by that of $N$ outside $S'$.
At $S'$ we restore the factors in $B_D$ using the indices of $M$.
These are $(e,f)=(8,4)$ at $2$, $(2,2)$ at $3$ and $5$, $(2,4)$ at
$7,11,13$, and $f=4$ at $17,19,23,29,31$. To justify these indices,
Proposition~\ref{ar:central-fields} identifies $M$ with the field of
$G/D_3G$, which retains every local cut of Section~\ref{tw:tower}.
A direct local computation from the radicands also confirms them.
They give the terms in \eqref{an:bad-restored} and the local floors
used below.

The conditions defining $\mathcal P_4$ and the census complement are
the tests \eqref{ar:nonsplit-test} on the central quadratic forms and
the twelve full-field predicates. These are tests in $M$. Hence a prime
in $\mathcal P_4$ has $f_M(p)=4$, and a complement prime has $f_M(p)=2$.
Both savings $S_4$ and $S_{\rm cen}$ are thus realized in $M$.
Each term in the following absolutely convergent series decreases in
$e_M(p)$ and $f_M(p)$, so the same local floors give
\[
 -\frac{\zeta_M'}{\zeta_M}(2)\cdot\frac1{2^{19}}
 =\sum_p\frac{\log p}{e_M(p)\,(p^{2f_M(p)}-1)}\leq R
\]
with $R$ as in \eqref{an:R}.
Consequently the fixed-field quantity
\[
 H_M:=\frac{\log\zeta_M(\sigma)}{2^{19}}
 +(\sigma-1)\left(\frac{\ell-\gamma_E-\log(4\pi)}4
 -\frac{\zeta_M'}{\zeta_M}(2)\cdot\frac1{2^{19}}\right)
\]
satisfies, with $L$ as in \eqref{an:slope} and the five allowances
\eqref{an:Y-numerical}, \eqref{an:bad-value},
\eqref{an:order-four-saving}, \eqref{an:census-value} and
\eqref{an:debit},
\[
 H_M\leq Y(\sigma)+B_D(\sigma)-S_4-S_{\rm cen}+(\sigma-1)L
 <0.04216181893948094953138458.
\]
The accompanying Lean formalization assumes this inequality for $H_M$
with the relaxed threshold $0.042165819$. Under that hypothesis, it
proves the transfer from $M$ to the tower fields and
Proposition~\ref{an:prime-budget}. The numerical bound for $H_M$ itself
is not proved in the formalization.
\end{remark}

\subsection{The Numerical Prime-Budget Bound}

We specify the local floors used in Proposition~\ref{an:prime-budget}.
For $p\leq10^4$, take $e_2^0=8$, $e_p^0=2$ at
$p=3,5,7,11,13$, and $e_p^0=1$ elsewhere. Use the exact retained
residue degrees $4$ at $2$, $2$ at $3,5$, and $4$ at
$7,11,13,17,19,23,29,31$. At any other such prime use $f_p^0=4$
if its square is nonzero in the full central quotient, $f_p^0=2$
if its genus vector is nonzero, and $f_p^0=1$ otherwise.
We apply the square test only when the genus vector is nonzero. For a
genus-split prime with a nontrivial central image, we keep the weaker
floor one. All these are lower bounds supplied by the
actual finite quotients in Proposition~\ref{ar:central-fields}.
Above $10^4$ use both floors one.

Since $\log x/(x^2-1)$ decreases on $[2,\infty)$, the integer tail
is bounded by its integral. It is enough to take
\begin{equation}\label{an:R}
 R=\sum_{p\leq10^4}\frac{\log p}{e_p^0(p^{2f_p^0}-1)}
          +\frac{\log10^4+1}{10^4(1-10^{-8})}.
\end{equation}
All finite floors are reconstructed by exact genus and square tests.
Outward evaluation gives
\begin{align}
 B_D(\sigma)&<0.041727023150898786,\label{an:bad-value}\\
 L:=\frac{\ell-\gamma_E-\log(4\pi)}4+R
 &<0.824453118933538946712643997041,\label{an:slope}\\
 (\sigma-1)L&<0.00006870442657779491222606.
 \label{an:debit}
\end{align}
The slope is independent of the abscissa, and the debit uses exactly
$\sigma-1=1/12000$.

Combining Proposition~\ref{an:prime-budget} with
\eqref{an:full-euler}, and then using the five conservative rational
bounds \eqref{an:Y-numerical}, \eqref{an:bad-value},
\eqref{an:order-four-saving}, \eqref{an:census-value}, and
\eqref{an:debit}, gives
\begin{align*}
 \limsup_{d\to\infty}\frac{\log L_F(1)}d
 &\leq Y(\sigma)+B_D(\sigma)-S_4-S_{\rm cen}+(\sigma-1)L\\
 &<0.04216181893948094953138458\\
 &<0.042161819.
\end{align*}
The middle number is the exact sum of the five rational bounds, with
the saving terms subtracted using their lower endpoints. The eventual
uniformity in Proposition~\ref{an:prime-budget} now proves
Theorem~\ref{an:ceiling}.
